%=============================================================================!
% 15.8  DOES THE RUN SAMPLE ITS ENSEMBLE?                                     !
%=============================================================================!
\section{Does the run sample its ensemble?}
\label{sec:cv-ensemble}

An error bar says how well a run knows its own averages. It is silent on
whether they belong to the ensemble the run claims: Chapter~13 found
thermostats
that hold the mean temperature and get the spread of the energy wrong
(\cref{sec:th-berendsen}). This section derives a test that needs no
knowledge of the system, runs at two temperatures and the canonical
density alone\cite{Shirts2013}, and sets out the resampling that gives its
error bar.

\subsection*{The ratio of two densities}

In the canonical ensemble the probability of a microstate of energy $E$
is $\mathrm{e}^{-\beta E}/Z$ (\cref{sec:sm-canonical}). The microstates
with energies between $E$ and $E + \delta E$ number $\Omega(E)$, the
volume of that shell (\cref{sec:sm-micro}), and grouping them, as
\cref{sec:sm-kinetic} grouped the velocities by their kinetic energy,
gives the density of $E$ at the inverse temperature $\beta$, written
$\mathcal{P}(E \mid \beta)$:
\begin{equation*}
   \mathcal{P}(E \mid \beta) = \frac{\Omega(E)\,\mathrm{e}^{-\beta E}}
   {Z(\beta)\,\delta E} .
\end{equation*}
No run knows $\Omega(E)$. At two temperatures the ratio of the densities
removes it, with $\delta E$:
\begin{equation*}
   \frac{\mathcal{P}(E \mid \beta_2)}{\mathcal{P}(E \mid \beta_1)}
   = \frac{Z(\beta_1)}{Z(\beta_2)}\,\mathrm{e}^{-(\beta_2 - \beta_1)E} ,
   \qquad
   \ln\frac{\mathcal{P}(E \mid \beta_2)}{\mathcal{P}(E \mid \beta_1)}
   = \ln\frac{Z(\beta_1)}{Z(\beta_2)} + (\beta_1 - \beta_2)E .
\end{equation*}
The logarithm of the ratio of the two histograms of the energy must be a
straight line in $E$ whose slope is $\beta_1 - \beta_2$, known in advance,
whatever the system. The intercept, which holds the unknown $Z$, does not
matter. For 130 and \SI{140}{\kelvin} the slope is
\SI{6.376}{\per\electronvolt}. The test needs the two histograms to
overlap. With $C_V = 2.347\,N\kB$ for the liquid (\cref{sec:sm-kinetic}),
the mean energy rises by about $C_V\Delta T = \SI{0.518}{\electronvolt}$
(\SI{0.520}{\electronvolt} between the runs below), about twice the
canonical spread $\sqrt{\kB T^2C_V}$, \SI{0.275}{\electronvolt} at
\SI{130}{\kelvin} and \SI{0.296}{\electronvolt} at \SI{140}{\kelvin}
(\cref{eq:sm-fluctuation}), so the two share a range of energies. On
100\,000 independent energies at each temperature drawn from a gamma
density (\cref{sec:sm-kinetic}), which has the form $\Omega(E)\mathrm{e}^{-\beta
E}$ with $\Omega \propto E^{399}$, a shape chosen only so that the two
densities overlap, the fitted slope is $(6.350 \pm
0.030)$\,\si{\per\electronvolt}, its error from the bootstrap below, 0.9
errors from the required value.

\subsection*{The bootstrap}

The slope is fitted to histograms, and its error is not that of a mean.
Efron's bootstrap finds the error of any quantity computed from samples
by recomputing it on new sets, each drawn from the samples themselves
with replacement and each as large as the original set, and taking the
spread of the results\cite{Efron1979}. Drawing single samples would destroy the memory of
a run and give too small an error; drawing whole blocks, longer than the
memory, keeps it within each block, as in \cref{sec:cv-blocks}, and makes
the resampled runs nearly as correlated as the real one, losing the memory
only at the joins between blocks\cite{Kunsch1989} (\texttt{mdlab} draws
non-overlapping blocks, a simpler variant of Künsch's moving blocks). The
function \texttt{block\_bootstrap} of \texttt{mdlab.analysis.stats} does
this for several runs at once, resampling each independently. On the
\SI{2}{\nano\second} run, blocks of \SI{20}{\pico\second} give the error
of the mean potential energy as \SI{0.0045}{\electronvolt}, against
\SI{0.0047}{\electronvolt} from $g$. \Cref{sec:cv-observables} used the
same resampling for the error of a heat capacity, a variance rather than a
mean.

\subsection*{Two thermostats}

\Cref{fig:cv-ensemble} runs the liquid for \SI{1}{\nano\second} at 130 and
\SI{140}{\kelvin} under CSVR with $\tau_{\mathrm T} = \SI{1}{\pico\second}$,
and under Berendsen coupling with $\tau_{\mathrm T} = \SI{0.1}{\pico
\second}$. After the start found from each run's total energy, the mean
energies agree within 0.7 of their combined errors: $-8.6409$ and
$-8.6306$\,\si{\electronvolt} at \SI{130}{\kelvin}, $-8.1209$ and
$-8.1118$\,\si{\electronvolt} at \SI{140}{\kelvin}. The spreads do not:
0.273 and \SI{0.287}{\electronvolt} under CSVR, close to the canonical 0.275
and \SI{0.296}{\electronvolt}, and 0.112 and
\SI{0.118}{\electronvolt} under Berendsen coupling. The slope of the
logarithm of the ratio, with its error from the block bootstrap with
blocks of \SI{20}{\pico\second}, is $(6.72 \pm 0.41)$\,\si{\per\electronvolt}
under CSVR, $1.05 \pm 0.06$ times the required value and 0.8 errors from
it: CSVR passes. Under Berendsen coupling it is $(35.1 \pm
3.5)$\,\si{\per\electronvolt}, 5.5 times the required value and 8.2 errors
from it. For two Gaussian densities of the same spread $\sigma_E$ whose
means differ by $\Delta\bar E$, the logarithm of the ratio, $[(E - \bar
E_1)^2 - (E - \bar E_2)^2]/2\sigma_E^2$, is a straight line of slope
$\Delta\bar E/\sigma_E^2$. The canonical $\Delta\bar E = C_V\Delta T$ and
$\sigma_E^2 = \kB T^2C_V$ make it $\Delta T/\kB T^2 \approx \beta_1 -
\beta_2$; Berendsen coupling keeps $\Delta\bar E$ and narrows $\sigma_E$ by
a factor of 2.43, which makes the line $2.43^2 = 5.9$ times too steep, close
to the 5.5 found. A thermostat that gives the right mean and the
wrong spread fails this test at once, whatever quantity was to be
measured; the physical-validation package of Merz and Shirts carries out
this test and others for the common simulation codes\cite{Merz2018}.

\begin{figure}[htb]
\centering
\includegraphics{ch15_convergence/ensemble}
\caption{Does the run sample its ensemble? The liquid for
\SI{1}{\nano\second} at 130 (light) and \SI{140}{\kelvin} (dark), under
CSVR (teal) and Berendsen coupling with $\tau_{\mathrm T} =
\SI{0.1}{\pico\second}$ (ochre). (a) The densities of the total energy.
(b) The logarithm of the ratio of the densities at the two temperatures,
against the slope $\beta_{130} - \beta_{140}$ that the canonical ensemble
requires (dashed).}
\label{fig:cv-ensemble}
\end{figure}

\begin{keyidea}
In the canonical ensemble $\ln[\mathcal{P}(E \mid \beta_2)/\mathcal{P}(E
\mid \beta_1)]$ is a straight line in $E$ with the slope $\beta_1 -
\beta_2$, whatever the system; runs at two nearby temperatures test it.
The error of a fitted slope, or of any quantity that is not a mean, comes
from the block bootstrap, which resamples blocks longer than the memory,
recomputes the quantity and takes the spread of the results.
\end{keyidea}

\notebook{15}{test a thermostat at two temperatures}

\begin{selfcheck}
Why does the test not need $\Omega(E)$, and what would a run at the wrong
temperature, \SI{132}{\kelvin} where \SI{130}{\kelvin} was claimed, do to
the slope?
\end{selfcheck}
